% from ega4-16-fr.tex
\oldpage[IV-4]{50}
est exacte; en outre, pour tout $x\in X$, il existe un voisinage ouvert $U$
de $x$ tel que les restrictions à $U$ des homomorphismes de (16.9.13.1)
forment une suite exacte et scindée.
\end{proposition}

Prouvons d'abord le lemme suivant:

\begin{lemma}[16.9.13.2]
\label{IV.16.9.13.2}
Soient $A$ un anneau, $\mathfrak{I}$ un idéal de $A$,
$A'=A/\mathfrak{I}$, $(f_i)_{1\leq i\leq r}$ une suite d'éléments de $A$ qui
est $A'$-régulière, $\mathfrak{K}=\sum_i f_iA$,
$\mathfrak{L}=\mathfrak{I}+\mathfrak{K}$,
$\mathfrak{K}'=\sum_i f_iA'$, de sorte que $C=A/\mathfrak{L}$ est isomorphe à
$A'/\mathfrak{K}'$. Alors pour tout entier $n>0$ et tout entier $N\geq n$,
on a la relation
\[
  \mathfrak{I}\cap\mathfrak{K}^n
  =\mathfrak{I}\mathfrak{K}^n+(\mathfrak{I}\cap\mathfrak{K}^N).
  \tag{16.9.13.3}
  \label{IV.16.9.13.3}
\]
\end{lemma}

Il suffit évidemment de prouver que tout élément du premier membre est
contenu dans le second, et, par récurrence sur $n$, on se ramène au cas où
$N=n+1$. Un élément du premier membre de (16.9.13.3), étant dans
$\mathfrak{K}^n$, s'écrit $P(f_1,\ldots,f_r)$, où
$P\in A[T_1,\ldots,T_r]$ est homogène et de degré $n$. Si $f_i'$ est l'image
canonique de $f_i$ dans $A'$, l'hypothèse
$P(f_1,\ldots,f_r)\in\mathfrak{I}$ signifie que
$P(f_1',\ldots,f_r')=0$. Mais
$P(f_1',\ldots,f_r')\in\mathfrak{K}'^n$, donc l'image canonique de
$P(f_1',\ldots,f_r')$ dans
$\mathfrak{K}'^n/\mathfrak{K}'^{n+1}$ est nulle. Or l'hypothèse que la suite
$(f_i)$ est $A'$-régulière implique que l'homomorphisme canonique
\[
  \mathbf{S}_C^n(\mathfrak{K}'/\mathfrak{K}'^2)
  \longrightarrow\mathfrak{K}'^n/\mathfrak{K}'^{n+1}
\]
est bijectif (0, 15.1.9); on conclut que les coefficients de $P$ appartiennent
à $\mathfrak{L}=\mathfrak{I}+\mathfrak{K}$. Il en résulte aussitôt que l'on a
$P(f_1,\ldots,f_r)\in\mathfrak{I}\mathfrak{K}^n+\mathfrak{K}^{n+1}$, et comme
$P(f_1,\ldots,f_r)\in\mathfrak{I}$, on a finalement
\[
  P(f_1,\ldots,f_r)
  \in\mathfrak{I}\mathfrak{K}^n
    +(\mathfrak{I}\cap\mathfrak{K}^{n+1}),
\]
ce qui prouve le lemme.

En prenant les quotients des deux membres de (16.9.13.3) par
$\mathfrak{I}\mathfrak{K}^n$, on voit que les relations (16.9.13.3) pour
$N\geq n$ entraînent
\[
  (\mathfrak{I}\cap\mathfrak{K}^n)/\mathfrak{I}\mathfrak{K}^n
  \subset\bigcap_{N\geq n}\mathfrak{K}^N
    \cdot(A/(\mathfrak{I}\mathfrak{K}^n)).
  \tag{16.9.13.4}
  \label{IV.16.9.13.4}
\]
On en déduit le

\begin{corollary}[16.9.13.5]
\label{IV.16.9.13.5}
Supposons vérifiées les hypothèses de (16.9.13.2) et supposons en outre que
l'anneau $A$ soit noethérien et que $\mathfrak{K}$ soit contenu dans le
radical de $A$. Alors on a pour tout entier $n>0$,
\[
  \mathfrak{I}\cap\mathfrak{K}^n=\mathfrak{I}\mathfrak{K}^n.
  \tag{16.9.13.6}
  \label{IV.16.9.13.6}
\]
\end{corollary}

En effet, le second membre de (16.9.13.4) est alors nul, puisque
$A/\mathfrak{I}\mathfrak{K}^n$ est un $A$-module de type fini (Bourbaki,
\emph{Alg. comm.}, chap.~III, §~3, n\textsuperscript{o}~3, prop.~6).

Prenons en particulier $n=2$ dans (16.9.13.6), et remarquons que l'on a
$\mathfrak{L}^2=\mathfrak{I}^2+\mathfrak{I}\mathfrak{K}+\mathfrak{K}^2
=\mathfrak{I}\mathfrak{L}+\mathfrak{K}^2$; puisque
$\mathfrak{I}\mathfrak{L}\subset\mathfrak{L}^2$, on en déduit que
\[
  \mathfrak{I}\cap\mathfrak{L}^2
  =\mathfrak{I}\mathfrak{L}+(\mathfrak{I}\cap\mathfrak{K}^2)
  =\mathfrak{I}\mathfrak{L}+\mathfrak{I}\mathfrak{K}^2
  =\mathfrak{I}\mathfrak{L},
\]
autrement dit
\[
  \mathfrak{I}\cap\mathfrak{L}^2=\mathfrak{I}\mathfrak{L},
  \tag{16.9.13.7}
  \label{IV.16.9.13.7}
\]
ce qu'on peut encore exprimer en disant que l'homomorphisme canonique
\[
  \mathfrak{I}/\mathfrak{I}\mathfrak{L}
  \longrightarrow(\mathfrak{I}+\mathfrak{L}^2)/\mathfrak{L}^2
\]
est bijectif.

Ces lemmes étant démontrés, prouvons la première assertion de (16.9.13): il
suffit
